\DocumentMetadata{
  pdfversion=2.0,
  pdfstandard=ua-2,
  lang=en-US,
  tagging=on,
  tagging-setup={math/setup=mathml-SE,tabsorder=structure}
}
\documentclass[11pt,letterpaper]{article}
\usepackage[margin=0.68in,includefoot]{geometry}
\usepackage{newcomputermodern}
\usepackage{amsmath,amscd,mathtools}
\providecommand{\square}{\mdlgwhtsquare}
\usepackage[hidelinks]{hyperref}
\hypersetup{
  pdftitle={Probability PhD Exam, September 1997},
  pdfauthor={University of Florida Department of Mathematics},
  pdfsubject={Graduate mathematics examination},
  pdfdisplaydoctitle=true,
  bookmarksopen=true
}
\setcounter{secnumdepth}{0}
\setlength{\parindent}{0pt}
\setlength{\parskip}{0.28em}
\setlength{\emergencystretch}{3em}
\setlength{\leftmargini}{2.15em}
\setlength{\leftmarginii}{2.65em}
\setlength{\leftmarginiii}{3em}
\renewcommand{\labelenumi}{\textbf{\theenumi.}}
\pagestyle{empty}
\raggedbottom
\allowdisplaybreaks
\newenvironment{examproblems}[1]{%
  \begin{enumerate}%
  \setcounter{enumi}{#1}%
  \setlength{\topsep}{0.35em}%
  \setlength{\partopsep}{0pt}%
  \setlength{\itemsep}{0.48em}%
  \setlength{\parsep}{0.18em}%
}{\end{enumerate}}
\newenvironment{examsubparts}{%
  \begin{enumerate}%
  \setlength{\topsep}{0.18em}%
  \setlength{\partopsep}{0pt}%
  \setlength{\itemsep}{0.16em}%
  \setlength{\parsep}{0.1em}%
}{\end{enumerate}}
\newenvironment{examinstructions}{%
  \par\begingroup\itshape
}{%
  \par\endgroup\vspace{0.18em}
}
\makeatletter
\renewcommand\section{\@startsection{section}{1}{0pt}%
  {-1.25ex plus -.25ex minus -.1ex}%
  {0.55ex plus .1ex}%
  {\normalfont\large\bfseries}}
\renewcommand{\@maketitle}{%
  \newpage\null\vskip -1.2em%
  {\footnotesize\bfseries
    \href{https://www.ufl.edu/}{University of Florida}, \href{https://math.ufl.edu/}{Department of Mathematics}\par}%
  \vskip 0.18em%
  \tagstructbegin{tag=Title}%
    {\LARGE\bfseries\href{https://gma.math.ufl.edu/exams/probability-phd/}{Probability PhD Exam}, September 1997\par}%
  \tagstructend%
  \vskip 0.35em%
  \tagmcbegin{artifact}%
    \hrule height 1pt%
  \tagmcend%
  \vskip 0.55em%
}
\makeatother
\title{Probability PhD Exam, September 1997}
\author{}
\date{}
\begin{document}
\maketitle
\thispagestyle{empty}
\begin{examproblems}{0}
\item
Let~\(X\) and~\(Y\) be independent and uniform on \([0,1]\). Find
\begin{examsubparts}
\item[\textnormal{a)}]
the distribution of \(X+Y\).
\item[\textnormal{b)}]
the conditional density of~\(X\) given \(X+Y=z\).
\end{examsubparts}
\item
Let \(X_0,X_1,\ldots,X_n,\ldots\) be i.i.d. with mean~\(m\). Let~\(N\) be Poisson with mean~\(\lambda\), independent of the~\(X\)’s. Define 
\[
Y=X_0+\cdots+X_N.
\]
 Prove that~\(Y\) is integrable. Find \(E(Y)\).
\item
State precisely (all for i.i.d.):
\begin{examsubparts}
\item[\textnormal{a)}]
the strong law of large numbers and prove it.
\item[\textnormal{b)}]
the central limit theorem.
\item[\textnormal{c)}]
the law of the iterated logarithm.
\end{examsubparts}
\item
Define a martingale relative to a filtration \(\{\mathcal F_n\}\), \(n=0,1,2,\ldots\). Define a stopping time relative to this filtration. State the optional sampling theorem. State and prove the martingale convergence theorem.
\item
What is an infinitely divisible distribution on \(\mathbb R^1\)? Prove that a non-trivial infinitely divisible distribution cannot have compact support.
\item
Define a standard Brownian Motion~\(W\) starting at 0. Define new processes \(W_1,W_2,W_3\) by: 
\[
W_1(t)=cW\!\left(\frac{t}{c^2}\right),\qquad t\ge 0,\quad c\text{ real}\ne0,
\]
 
\[
W_2(t)=tW\!\left(\frac1t\right),\qquad t>0,\quad W_2(0)=0,
\]
 
\[
W_3(t)=
\begin{cases}
W(1)-W(1-t),&0\le t\le1,\\
W(t),&\text{otherwise}.
\end{cases}
\]
 Show that \(W_1,W_2,W_3\) are the standard Brownian motions.
\item
Let \(\varphi(t)=\int e^{itx}\,\mu(dx)\) be a characteristic function where~\(\mu\) is a probability measure on \(\mathbb R^1\).
\begin{examsubparts}
\item[\textnormal{a)}]
Suppose \(|\varphi(t)|=1\) for some \(t\ne0\). Then prove that unless \(|\varphi(t)|\equiv1\), there is a smallest \(t_0\ne0\) and a~\(d\) such that~\(\mu\) is concentrated on the set 
\[
\left\{d+\frac{2\pi j}{t_0}\right\},\qquad j=0,\pm1,\pm2,\ldots.
\]
\item[\textnormal{b)}]
Using a) prove that \(|\cos t|\) is not a characteristic function. Hint: \(\cos t\) and \(\cos^2t\) are characteristic functions.
\end{examsubparts}
\item
Let \(S_n=(U_n,V_n)\) denote the position after~\(n\) steps of a random walk on \(\mathbb Z^2\), starting from \((0,0)\). Suppose that 
\[
U_{n+1}=U_n\pm1,
\qquad
V_{n+1}=V_n\pm1,
\]
 where the signs are picked by two independent tosses of a fair coin, independently at each step. Define 
\[
p_n=P(S_n=(0,0)),
\]
 
\[
\tau_n=\inf\{m>\tau_{n-1}:S_m=(0,0)\}.
\]
 An easy fact is \(P(\tau_n<\infty)=P(\tau_1<\infty)^n\).
\begin{examsubparts}
\item[\textnormal{(1)}]
Prove that for any random walk \(S_n\in\mathbb Z^2\) the following are equivalent:
\begin{examsubparts}
\item[\textnormal{(i)}]
\(P(\tau_n<\infty)=1\).
\item[\textnormal{(ii)}]
\(P(S_n=(0,0)\text{ i.o.})=1\).
\item[\textnormal{(iii)}]
\(\displaystyle\sum_{n=1}^{\infty}P(S_n=(0,0))=\infty\).
\end{examsubparts}
\item[\textnormal{(2)}]
Find a formula for \(p_n\), \(n>0\).
\item[\textnormal{(3)}]
Use (1) and Stirling’s formula 
\[
n!\approx\sqrt{2\pi n}\left(\frac ne\right)^n e^S
\qquad(S\to0\text{ as }n\to\infty)
\]
 to show that the diagonal random walk in \(\mathbb Z^2\) is recurrent.
\end{examsubparts}
\item
Let \((P_t)_{t>0}\) be a Markov semigroup. Assume \(P_t(x,dy)=p_t(x,y)\,dy\) with \(p_t(x,y)=p_t(y,x)\) and \(p_t(x,y)\le M(t)\), where \(M(t)\) is a constant depending on~\(t\). If~\(\mu\) is a finite signed measure, show that 
\[
\int p_t(x,y)\,\mu(dx)\mu(dy)\ge0.
\]
\item
An urn at \(t=0\) contains \(R_0\) (\(\ge1\)) red balls and \(B_0\) (\(\ge1\)) black balls, and “random sampling with double replacement” is called out as follows:

\par
(i) Choose a ball at random (meaning that each ball in the urn has the same probability to be drawn), and note its color;

\par
(ii) Replace this ball in the urn together with an extra ball of the same color;

\par
(iii) Go to (i).

\par
Let \(M_t\) be the proportion in red balls in the urn after~\(t\) such operations, so that 
\[
M_0=\frac{R_0}{R_0+B_0}.
\]
 Show that, with respect to a suitable filtration (which you should specify), the sequence \(M_0,M_1,\ldots,M_t,\ldots\) is a martingale. Prove that \(M_t\) converges as \(t\to+\infty\) almost surely to a random variable \(M_\infty\), \(0\le M_\infty\le1\). Show that 
\[
M_0=EM_1=EM_2=\cdots=EM_t=EM_\infty,
\]
 and that 
\[
M_0^2\le E[M_1^2]\le\cdots\le E[M_t^2]\le E[M_\infty^2],
\]
 and that 
\[
P(M_\infty=M_0)<1.
\]
\end{examproblems}
\begin{examinstructions}
NOTE. You may quote without proof any martingale theorems you need, but should state them carefully.
\end{examinstructions}
\begin{examproblems}{10}
\item
Let \(X_n\) be independent and all distributed according to standard normal distribution. Find constants \(C_n\) such that 
\[
\limsup \frac{X_n}{C_n}=1\qquad\text{a.s.}
\]
 Hint. 
\[
\int_a^\infty e^{-b^2/2}\,db\le\frac1a e^{-a^2/2},
\]
 and let \(c_n^2=2\log(n\log n)\).
\end{examproblems}
\end{document}
